\documentclass[11pt]{article}

\usepackage[utf8]{inputenc}
\usepackage[T1]{fontenc}
\usepackage{amsmath,amssymb,amsthm}
\usepackage{hyperref}
\usepackage[margin=1in]{geometry}

\theoremstyle{definition}
\newtheorem{definition}{Definition}[section]
\newtheorem{theorem}{Theorem}[section]
\newtheorem{lemma}[theorem]{Lemma}

\title{A Formal Verification of TCP (RFC~793) in Lean~4}
\author{Samuel Schlesinger}
\date{\today}

\begin{document}

\maketitle

\begin{abstract}
We present a machine-checked formalization of the Transmission Control
Protocol as specified in RFC~793. Using the Lean~4 theorem prover, we model
the TCP state machine, define segment acceptability predicates over modular
32-bit sequence number arithmetic, and prove correctness properties of
connection establishment (three-way handshake), graceful teardown, and global
safety invariants. All theorems are verified by Lean's kernel---if the code
typechecks, the proofs hold.
\end{abstract}

\tableofcontents

\section{Introduction}
\label{sec:introduction}

This work presents a machine-checked formalization of the Transmission Control
Protocol (TCP) as originally specified by Postel~\cite{rfc793}.

\section{Sequence Numbers}
\label{sec:seqnum}

% Modular 32-bit arithmetic, the SeqNum type, and comparison relations.
% Corresponds to \texttt{Tcp.SeqNum}.

\section{Core Types}
\label{sec:types}

% Port, Socket, Segment header, and TCB structures.
% Corresponds to \texttt{Tcp.Types}.

\section{State Machine}
\label{sec:state}

% The eleven TCP states, event types, and the transition function.
% Corresponds to \texttt{Tcp.State}.

\section{Segment Acceptability}
\label{sec:acceptability}

% Acceptability predicates for segments and ACKs.
% Corresponds to \texttt{Tcp.Acceptability}.

\section{Connection Establishment}
\label{sec:handshake}

% Three-way handshake correctness theorems.
% Corresponds to \texttt{Tcp.Handshake}.

\section{Connection Teardown}
\label{sec:teardown}

% Graceful close and TIME-WAIT properties.
% Corresponds to \texttt{Tcp.Teardown}.

\section{Invariants and Safety Properties}
\label{sec:invariants}

% Global safety invariants and data delivery guarantees.
% Corresponds to \texttt{Tcp.Invariants}.

\section{Conclusion}
\label{sec:conclusion}

% Summary of results, limitations, and future work.

\bibliographystyle{plain}
\bibliography{references}

\end{document}
